% Supporting inherited certificate only; not an input to the main theory.
\section{An inherited compact certificate as a supporting example}
\label{app:compact-example}

This appendix applies the analytical screen of
Proposition~\ref{cmp:small-slice} to an existing certificate. The source
values are numerical supporting evidence, already obtained with complete
source, interpolation, quadrature and arithmetic allowances in
\cite{BakerGT15,BakerGT14}; no new source evaluation is performed here.
The main contour, lattice and comparison theorems do not depend on them.
The resulting arithmetic and scope were checked in \cite{BakerCM27}.

The certificate follows the Lehmer pair near $a=14010.1$ on
$0\le t\le0.1$. The gauged field $G=\e^{\pi a/8}H$ counts two zeros
inside $|z-a|<3/4$ and has a cleared annulus out to radius $21/20$.
Its centered annular logarithm has modulus less than $1.737$ on both
boundaries and hence throughout the annulus. Its directly enclosed
incoming values at $t_*=0.1$ imply the safe rounded bounds
\begin{equation}
 \begin{gathered}
 0.1787368\le D(t_*)\le0.1787369,\qquad
 0.6037742\le M(t_*)\le0.6037744,\\
 V(0)\le1-2(0.1787368)(0.6037742)
       =0.78416666313888.
 \end{gathered}
 \label{supp:incoming}
\end{equation}
The older certificate's drift symbol is $K_{\mathrm{GT}}=-M$; it is
not the forcing $K=2D^2\mathcal F$ used in the present paper.
Here $V(0)$ is at backward time zero, namely physical time $t_*$.

Consider the proposed local descent to $t=-0.0003$, of duration
$\Delta=0.1003$. With $d_{\mathrm{stop}}=0.0001$, a hypothetical
$v\le0.89$ would leave the strict margin
\[
 0.1787368-2(0.89)(0.1003)-0.0001=0.0001028>0,
\]
and would require the forcing floor $\kappa\ge1-0.89^2=0.2079$.
The existing certificate ends at physical time zero: neither the
negative-time source bounds nor the proposed floor are available from it.
Even the favorable initial-$M$ scalar allowance is at most
\[
 \frac{0.1787369}{2(1-0.1787369\cdot0.6037744)}<0.100180,
\]
which does not reach the proposed duration.

The already certified physical slice $t=0$ is enough to test the stronger
constant-$v$ proposal. At that slice,
\[
 |c-a|<0.1,\qquad 0<D<0.002,\qquad
 |\mathcal L-\mathcal L_{\mathrm{center}}|<7/4
 \quad\text{on }|z-a|=0.9.
\]
There $|Q|>0.638>3/5$ and $|Q'|\le2$. The complete log kernel
\eqref{cmp:log-kernel} gives
\begin{align*}
 |S_2|&\le 2(0.9)(7/4)2/(3/5)^3=175/6,\\
 |S_3|&\le 3(0.9)(7/4)2/(3/5)^4=875/12,\\
 |\mathcal F|&\le12|S_2|+16(0.002)|S_3|\le1057/3.
\end{align*}
Consequently the actual complete forcing satisfies
\begin{equation}
 |K(t=0)|<2(0.002)^2(1057/3)
        =1057/375000<0.003.
 \label{supp:small-K}
\end{equation}
This upper bound is not a failed lower enclosure. Any constant floor
covering this slice has $\kappa<0.003$, and the selected constant
parameter must obey $v>\sqrt{0.997}>0.998$. Even exact initial data and
zero stop threshold therefore give only
\[
 \Delta<\frac{0.1787369}{2(0.998)}<0.0896<0.1003.
\]
Refining the numerical floor cannot make this particular consumer pass.
The slice is inside the proposed stopping region, since its certified
$D(t=0)>0.00142117>0.0001$. It is used to reject a new allowance, not to
present already known positivity as new coverage.

The conclusion concerns the one constant-$v$ barrier from this fixed
anchor over an interval containing physical time zero. It proves no
collision or failure of actual descent, and excludes neither a sharper
transient comparison nor time-dependent or regular coupled evolution.
The same-source direct endpoint estimate is stronger on the old interval;
no method is certified here on the proposed negative-time extension.
No global all-real start, practical all-height cutoff or improved bound
on $\Lambda$ follows from this local supporting example.

\subsection{An unsigned force envelope at the same anchor}
\label{supp:force-envelope}

The subsequent comparison in \cite{BakerCM32,BakerCMReview29} uses
the same anchor to test a specific compression of the common node-force
equation. For two nodes put $r_U=|PU|$, where $P$ subtracts the mean
and $U_i=H_{zz}(x_i)/(2H_z(x_i))$. The exact common-label identity gives
\[
 g(u)=g(0)-2\int_0^u(U_1-U_0)\dd s
       \ge g(0)-2\sqrt2\int_0^u r_U\dd s.
\]
In backward time the internal graph term is antidissipative. Replacing
its gap denominator by the fixed stop
$s=2\sqrt{d_{\rm stop}}=1/50$ and the exterior forcing by an
independent nonnegative norm allowance gives the particular scalar
envelope
\[
 \mathcal R'=10000\mathcal R+J_{\rm ext},\qquad
 J_{\rm ext}\ge0,\qquad \mathcal R(0)\ge r_U(0).
\]
Even with zero exterior allowance and zero added errors,
$\mathcal R(u)\ge r_U(0)\e^{10000u}$.
This is a property of the envelope, not a lower growth estimate for
the actual force norm.

The safe incoming bounds in \eqref{supp:incoming} imply
$g(0)<423/500$ and $M(0)<151/250$. The edge force identity therefore
gives
\[
 U_1-U_0=\frac2{g(0)}-g(0)M(0)
  >\frac{1000}{423}-\frac{423}{500}\frac{151}{250}
  >\frac{37}{20}.
\]
Since $2\sqrt2r_U(0)=2|U_1-U_0|>37/10$, a positive margin from this
unsigned envelope requires
\begin{equation}
 \Delta<10^{-4}\log\left(1+
       \frac{10000(0.846-0.02)}{3.7}\right)
       =10^{-4}\log\frac{82637}{37}<0.0008.
 \label{supp:force-duration}
\end{equation}
The last strict comparison follows exactly from
\[
 \e^8>\sum_{k=0}^{10}\frac{8^k}{k!}
       =\frac{3830591}{1575}>\frac{82637}{37}.
\]
This rejects the fixed-stop-denominator, unsigned-exterior scalar
compression for the proposed $0.1003$ descent. It does not reject
state-dependent denominators, signed exterior correlations, other blocks
or the complete force dynamics.

\subsection{The retained constant-\texorpdfstring{$K$}{K} transient still misses the endpoint}
\label{supp:constant-K-transient}

Theorem~\ref{rf:constant-K} retains more than the earlier constant-$v$
barrier. It can nevertheless be screened using the same actual included
slice \eqref{supp:small-K}. Give its scalar model the favorable data
\[
 d_+=0.1787369,\qquad b_+=0.6037744,\qquad
 v_-=1-2d_+b_+,\qquad \kappa_*=1057/375000.
\]
The initial gap $2\sqrt D$ and initial backward velocity
$-2/\sqrt D+4M\sqrt D$ are both increasing on the positive anchor
rectangle. The stated upper endpoints therefore maximize both initial
data. Increasing $\kappa$ likewise favors the cooperative comparison.
Any constant floor covering the certified $t=0$ slice is less than
$\kappa_*$, so these choices are favorable even to an exact-data use.

Substitution of the rational values into
\eqref{rf:constant-K-polynomial} gives the exact inequality
\begin{equation}
 d_+-2v_-\Delta+
       \frac{v_-^2+\kappa_*-1}{d_+}\Delta^2
       <-\frac{8239}{10^8}<0,
 \qquad \Delta=1003/10000.
 \label{supp:transient-failure}
\end{equation}
The model starts positive and has reached its first zero before this
endpoint. Its negative polynomial value is not asserted as a
positive-gap comparison after that zero. Equivalently, even the
favorable endpoint requirement $d_\kappa(\Delta)>d_{\rm stop}$ would
need
\[
 \kappa>
 1-v_-^2+\frac{d_+}{\Delta^2}
       (d_{\rm stop}-d_++2v_-\Delta)
 >0.00605917>\kappa_*.
\]
For a stopped-trial use, successful passage would put the actual included
slice inside the trial region. If a lower stop occurred earlier, the
proposed target would already have failed. The necessary screen therefore
does not assume the desired whole-interval positivity.

This strengthens the earlier allowance rejection to the sharp scalar
constant-$K$ transient comparison from this anchor. It still excludes
neither time-dependent $K$ nor a constant physical $\mathcal F$ floor.

\subsection{A positive conditional margin for retained physical forcing}
\label{supp:retained-F}

Set
\[
 d=0.1787368,\qquad b=0.6037742,\qquad f=3/10,\qquad
 \Delta=1003/10000.
\]
These are the incoming lower data and a proposed complete forcing floor,
not a measured lower value of $\mathcal F$. The model of
Theorem~\ref{rf:forcing-comparison} has a rigorously positive margin.
Here is the exact rational calculation, as recorded in
\cite{BakerCM32,BakerCMReview29}.

Put $\lambda^2=2f=3/5$, $w=4\lambda^2\Delta^2$, $v=1-2db>0$ and
\[
 J=\frac d2+\frac{2db^2-2b}{\lambda^2}<0.
\]
Expanding \eqref{rf:physical-forcing-model} by double-angle identities
gives
\begin{equation}
 d_f(\Delta)=d+J[\cosh(2\lambda\Delta)-1]
                    -v\frac{\sinh(2\lambda\Delta)}{\lambda}.
 \label{supp:F-expansion}
\end{equation}
The positive Taylor series have the rational upper bounds
\begin{align*}
 C_+&=\sum_{j=0}^4\frac{w^j}{(2j)!}
       +\frac{w^5/10!}{1-w/(12\cdot11)},\\
 S_+&=2\Delta\left[
        \sum_{j=0}^4\frac{w^j}{(2j+1)!}
        +\frac{w^5/11!}{1-w/(13\cdot12)}\right].
\end{align*}
Indeed, after the first omitted terms the successive ratios are at
most $w/(12\cdot11)$ and $w/(13\cdot12)$, respectively; both are
strictly below one. Thus
$\cosh(2\lambda\Delta)\le C_+$ and
$\sinh(2\lambda\Delta)/\lambda\le S_+$.
The two coefficients $J$ and $-v$ in \eqref{supp:F-expansion} are
negative, so exact rational substitution proves
\begin{equation}
 d_f(\Delta)\ge d+J(C_+-1)-vS_+
       >\frac{16246}{10^8},\qquad
 d_f(\Delta)-d_{\rm stop}>\frac{6246}{10^8}.
 \label{supp:F-margin}
\end{equation}
All quantities in this comparison are rational; no sampling of the
canonical source or fit to a time profile enters it. The inequalities
$b>\sqrt{f/2}$ and $4bd<2$ hold for these data. Corollary
\ref{rf:forcing-stop} therefore makes the endpoint the model's minimum
on the full interval, so the second inequality is the strict lower-stop
margin, conditional on the actual source hypotheses.

For the same trial geometry $r=9/10$, $C=1/10$, $D_{\max}=1/5$ and
$q=11/25$, \eqref{rf:source-errors} gives the exact amplification factors
\[
 E_{\mathcal F,B}\le\frac{119250}{1331}e_B,
 \qquad
 E_{\mathcal F,\mathcal L}\le\frac{14175000}{14641}e_{\mathcal L}.
\]
The complete source/time error, the signed lower floor
$\mathcal F_\Gamma\ge3/10$, the negative-time contour clearance and
the unconditional coordinate bounds are unacquired. The annular
logarithm bound used earlier provides an absolute forcing bound, not
this signed lower floor. Small $K=2D^2\mathcal F$ at $t=0$ is compatible
with this proposed physical floor, but does not prove it. Direct regular
$D/M$ evolution with identical data and errors gives precisely the
same model and margin. Thus \eqref{supp:F-margin} is a conditional
consumer calculation, not a new canonical certificate or a lattice
advantage.

\subsection{A distinct time-zero census and its forward constants}
\label{fwd:census-example}

The finite census used in \cite{BakerCM34,BakerCMReview30} lies at a
different height from the Lehmer pair. Its constants are not transferred
to that pair. The accepted source record \cite{BakerOCCensus22} encloses
every zero in
\[
 \mathcal R=(988,2012)+\ii(-3/2,3/2)
\]
in 388 occupied cells, each of count one including multiplicity.
Their horizontal projections are disjoint, their maximum horizontal
width is $5/32768$, and successive projections have separation at
least $5085/8192$. The common recorded imaginary interval is
$[-3/32768,3/131072]$, which is not symmetric about zero.

Reality follows from the complete census and conjugation, rather than
from replacing these imaginary intervals by zero. The conjugate of a
root in $\mathcal R$ is another root in the same symmetric outer
rectangle, with the same real coordinate. Completeness puts it in an
occupied cell; disjoint horizontal projections force that cell to be
the original one. Its count one forces equality of the two roots and
simple multiplicity. There are consequently 388 real simple nodes,
387 internal consecutive gaps and 385 edges with both neighboring gaps
inside the census. This deduction retains the inherited winding
certificate's scope; it is not an independent replay of its source
arithmetic.

The initial squared-half-gap minimum and the finite-block coefficient are
\begin{equation}
 q_0=\frac{5085^2}{4\cdot8192^2}
       =\frac{25857225}{268435456},\qquad
 a_{388}=\frac{388}{194^2}=\frac1{97}.
 \label{supp:census-constants}
\end{equation}
If a future common tube keeps the block diameter below $1024$, the
force-dispersion contraction coefficient $2N/L_{\max}^2$ is
$\gamma=97/131072$. Even with zero exterior norm forcing, its
homogeneous envelope over the illustrative time $s=1/10$ decreases by
a relative amount
\[
 1-\e^{-\gamma s}<\gamma s=\frac{97}{1310720},
\]
about $0.0074$ percent. This is the reduction in that conservative
envelope, not an upper bound on actual contraction. Continued
confinement and zero exterior forcing have not been acquired.

For a paid uniform upper exterior interaction $C\ge0$, the forward
minimum envelope is
\[
 q(s)\ge\e^{-4Cs}q_0+
          \frac{1-\e^{-4Cs}}{388C},
\]
with limit $q_0+s/97$ when $C=0$. Its lower bound exceeds $q_0$ exactly
when
\begin{equation}
 C<\frac{a_{388}}{4q_0}
   =\frac{268435456}{10032603300}.
 \label{supp:census-exterior-target}
\end{equation}
No such whole-time exterior bound is supplied by the census. Both block
ends must be included in its acquisition: the 12-unit buffers of the
smaller observation window $[1000,2000]$ cannot be assigned to the
outermost retained nodes. Likewise, selecting only interior eligible
edges does not establish the cumulative-distance inequalities for the
minimum over the complete block.

This finite artifact supplies actual time-zero labels and gap geometry,
not a positive-time source tube, an exterior-force bound, or an all-height
gap law. The analytical forward comparisons remain conditional on their
independent source and boundary payments, and Theorem~\ref{rf:horizon}
still applies to any subsequent use of its specific backward consumer.
